%==============================================================================
% Sjabloon poster bachproef
%==============================================================================
% Gebaseerd op document class `a0poster' door Gerlinde Kettl en Matthias Weiser
% Aangepast voor gebruik aan HOGENT door Jens Buysse en Bert Van Vreckem

\documentclass[a0,portrait]{hogent-poster}

% Info over de opleiding
\course{Bachelorproef}
\studyprogramme{toegepaste informatica}
\academicyear{2025-2026}
\institution{Hogeschool Gent, Valentin Vaerwyckweg 1, 9000 Gent}

% Info over de bachelorproef
\title{Digitale toegankelijkheid aan HOGENT vanuit het perspectief van studenten met een beperking.}
%\subtitle{Ondertitel (eventueel)}
\author{Eray Köksoy}
\email{eray.koksoy@student.hogent.be}
\supervisor{Lotte Van Steenberghe}
\cosupervisor{Alessandro De Waele (HOGENT)}

% Indien ingevuld, wordt deze informatie toegevoegd aan het einde van de
% abstract. Zet in commentaar als je dit niet wilt.
\specialisation{Full stack development}
\keywords{digitale toegankelijkheid, studenten met een beperking, HOGENT}
%\projectrepo{https://github.com/ErayKoksoy/latex-hogent-bachproef-Eray2026}

\begin{document}

\maketitle

\begin{abstract}
Digitale toegankelijkheid is een essentiële voorwaarde om gelijke participatiekansen in het hoger onderwijs te garanderen. De HOGENT-campusapp werd ontwikkeld als een proof of concept voor een centrale digitale toegangspoort tot studie-informatie, communicatie en planning. Wanneer toegankelijkheidsrichtlijnen onvoldoende worden toegepast, kunnen echter drempels ontstaan voor studenten met een beperking.

Deze bachelorproef onderzoekt hoe de HOGENT-campusapp inclusiever kan worden gemaakt door middel van een toegankelijkheidsaudit, de implementatie van verbeteringen volgens de WCAG-richtlijnen en een evaluatie van de impact op zowel toegankelijkheid als performantie. De resultaten tonen aan dat gerichte toegankelijkheidsverbeteringen leiden tot een hogere toegankelijkheidsscore zonder significante negatieve invloed op de performantie van de applicatie.
\end{abstract}

\begin{multicols}{2} % This is how many columns your poster will be broken into, a portrait poster is generally split into 2 columns

\section{Introductie}

Digitale toepassingen spelen een steeds belangrijkere rol binnen het hoger onderwijs. De HOGENT-campusapp werd ontwikkeld als een centrale digitale toegangspoort voor studenten, waarin functies zoals lesroosters, nieuwsberichten, deadlines, meldingen en campusinformatie worden samengebracht. Toch voldoet ongeveer 72\% van de Belgische overheidswebsites niet aan de geldende toegankelijkheidsnormen\textsuperscript{1}. Bovendien leeft 16,3\% van de jongeren tussen 16 en 24 jaar met een beperking\textsuperscript{2}, een groep die een aanzienlijk deel van de studentenpopulatie vertegenwoordigt. Ontoegankelijke digitale toepassingen kunnen daardoor leiden tot drempels bij het volgen van opleidingen en het raadplegen van essentiële informatie. Vanuit deze problematiek onderzoekt deze bachelorproef hoe de HOGENT-campusapp inclusiever kan worden ontworpen voor studenten met een beperking, rekening houdend met hun specifieke noden op het vlak van digitale toegankelijkheid.

\section{Experimenten}

Het onderzoek werd uitgevoerd in vijf opeenvolgende fasen:

\begin{enumerate}
    \item \textbf{Literatuurstudie:} WCAG 2.1 en 2.2, Europese en Belgische regelgeving, evaluatiemethoden voor digitale toegankelijkheid.

    \item \textbf{Toegankelijkheidsaudit:} Lighthouse, Accessibility Insights en handmatige toetsenbordtesten.

    \item \textbf{Probleemanalyse:} Identificatie van WCAG-overtredingen en prioritering op basis van gebruikersimpact.

    \item \textbf{Implementatie:} Verbetering van ARIA-labels, focusindicatoren, kleurcontrasten en toetsenbordnavigatie.

    \item \textbf{Hertesting:} Nieuwe Lighthouse-audit, vergelijking met de nulmeting en analyse van performantie-impact.
\end{enumerate}

\section{Belangrijkste verbeteringen}

\begin{itemize}
    \item \textbf{Waarnembaarheid:} toegankelijke namen, betere kleurcontrasten en schermlezerondersteuning.
    \item \textbf{Bedienbaarheid:} focusindicatoren, toetsenbordnavigatie en bereikbaarheid van interactieve componenten.
    \item \textbf{Begrijpelijkheid:} consistentere interactiepatronen en duidelijkere labels.
    \item \textbf{Robuustheid:} correct gebruik van ARIA-attributen en verbeterde compatibiliteit met assistieve technologieën.
\end{itemize}

\section{Resultaten}
\subsection{Toegankelijkheid}

De toegankelijkheidscores verbeterden op vrijwel alle onderzochte pagina's:

\begin{center}
\begin{tabular}{l c c}
\textbf{Pagina} & \textbf{Voor} & \textbf{Na} \\
\hline
Home & 86 & 96 \\
Nieuws & 89 & 100 \\
Campus & 90 & 100 \\
Lesrooster & 86 & 100 \\
Deadlines & 82 & 100 \\
Notificaties & 87 & 100 \\
\end{tabular}

\textit{Schaal: 0-49 (Slecht), 50-89 (Verbetering nodig), 90-100 (Goed)}
\end{center}

\subsection{Performantie}
De hertesting toont aan dat de implementatie van toegankelijkheidsmaatregelen geen significante negatieve impact heeft op de performantie van de applicatie.

\begin{itemize}
  \item FCP bleef stabiel
  \item LCP vertoonde slechts minimale verschillen
  \item Speed Index bleef vergelijkbaar
  \item CLS bleef ongewijzigd
  \item TBT wijzigde slechts beperkt
\end{itemize}

De resultaten bevestigen dat toegankelijkheid en performantie geen tegenstrijdige doelstellingen hoeven te zijn.

\begin{center}
  \captionsetup{type=figure}
  \includegraphics[width=0.33\linewidth]{focus-1-after.PNG}
  \captionof{figure}{Homepagina van de HOGENT-campusapp na implementatie van toegankelijkheidsverbeteringen}
\end{center}

\section{Conclusies}

Deze bachelorproef toont aan dat de toegankelijkheid van de HOGENT-campusapp aanzienlijk verbeterd kan worden door gerichte aanpassingen op component- en UI-niveau.

De uitgevoerde verbeteringen verhogen de bruikbaarheid voor studenten met een beperking, verbeteren de conformiteit met WCAG-richtlijnen en veroorzaken geen significante performantievermindering.

Digitale toegankelijkheid moet daarom niet worden beschouwd als een optionele uitbreiding, maar als een structureel kwaliteitscriterium binnen het ontwikkelproces van moderne applicaties.

\section{Toekomstig onderzoek}

\begin{itemize}
  \item Gebruikerstesten met studenten met een beperking
  \item Evaluatie op mobiele iOS- en Android-platformen
  \item Integratie van toegankelijkheidscontroles in CI/CD-pipelines
  \item Onderzoek naar WCAG 2.2 en toekomstige WCAG 3.0-richtlijnen
  \item Uitbreiding van het toegankelijkheidskader naar andere HOGENT-platformen
\end{itemize}

\section{Bronnen}

\small

\textsuperscript{1} FOD BOSA (2024)

\textsuperscript{2} Raad van de Europese Unie (2025)

\end{multicols}
\end{document}